\clearpage\SourceLeaf{87}%
\pdfbookmark[1]{Noch etwas über die Bestimmung des Osterfestes}{osterfest1807}
\begin{center}
{\large NOCH ETWAS ÜBER DIE BESTIMMUNG DES OSTERFESTES.\par}
\medskip\hrule\medskip
{\small Braunschweigisches Magazin. 1807. September 12.\par}
\medskip\hrule
\end{center}\medskip
Die cyclische Berechnung des Osterfestes, wovon einigemal in diesen Blättern
die Rede gewesen ist, obgleich an sich nicht schwer, beruhet doch auf Gründen,
die gerade nicht jedermann bekannt sind. Sie setzt ferner Hülfstafeln voraus,
die auch nicht in eines jeden Händen sind, und wie das Beispiel des Hrn.
Superint. \textsc{Helmuth} gezeigt hat, selbst in gedruckten Schriften fehlerhaft sein
können. Gleichwohl tritt nicht selten der Fall ein, wo man gern für ein künftiges
Jahr, oder für ein vergangenes, wofür man keinen Kalender zur Hand hat,
das Datum des Ostertages wissen möchte. Manchem wird daher die Kenntniss
eines besondern Verfahrens zur Bestimmung des Osterfestes nicht unwillkommen
sein, wobei man von den Gründen und Kunstwörtern der cyclischen Berechnungsart
gar nichts zu wissen braucht, keiner Hülfstafeln bedarf, und blos ein
Paar ganz einfache Rechnungsoperationen zu machen hat. Ich habe diese Methode
bereits vor sieben Jahren in des Hrn. \textsc{von Zach} \textit{Monatlicher Correspondenz}
(Augustheft 1800) mitgetheilt; da sie indess von den Lesern dieses Magazins wohl
nur wenigen bekannt geworden sein wird, so ist vielleicht diesem und jenem damit
gedient, hier das Nothwendigste von jenem Verfahren von neuem erklärt zu
finden. Ich begnüge mich indess, hier die Regeln zur Bestimmung des Osterfestes
im Gregorianischen Kalender für das gegenwärtige Jahrhundert, die nächst%
\SourceLeaf{88}%
künftigen und die vergangenen, auf eine auch dem Ungeübtesten fassliche Art
vorzutragen, wozu ich noch die Regeln für dieselbe Rechnung im Julianischen
Kalender beifügen werde. Die Erweiterung der Regeln für die spätern Jahrhunderte,
so wie ihren Zusammenhang mit der ursprünglichen Einrichtung unsers
Kalenders muss man am angeführten Orte nachlesen. [S. 75. d. B.]

Zur Bestimmung des Osterfestes für jedes beliebige Jahr des gegenwärtigen
Jahrhunderts von 1800 bis 1899 incl. hat man folgende Regeln zu beobachten:

$1^{0}$. Die Jahrzahl wird mit 19 dividirt; auf den Quotienten wird gar keine
Rücksicht genommen, sondern nur der übrig bleibende Rest bemerkt, den wir,
um ihn von denjenigen Resten zu unterscheiden, die sich bei den folgenden Operationen
ergeben werden, den \textit{ersten} Rest nennen wollen. Falls die Division
aufgeht, wird 0 als Rest angenommen, was auch bei den übrigen Divisionen
zu bemerken ist.

$2^{0}$. Dieser erste Rest wird mit 19 multiplicirt, zum Producte 23 hinzugefügt,
und diese Summe mit 30 dividirt; auch hier wird blos der Rest bemerkt,
den wir den \textit{zweiten} Rest nennen.

$3^{0}$. Die Jahrzahl wird mit 4 dividirt; der Rest soll der \textit{dritte} Rest heissen.

$4^{0}$. Man mache eine Summe aus der vierfachen Jahrzahl, dem sechsfachen
zweiten Reste, dem doppelten dritten und der Zahl 4; diese Summe mit 7 dividirt,
gibt uns den \textit{vierten} Rest.

$5^{0}$. Die Summe des zweiten und vierten Restes zeigt nunmehr an, wie
viele Tage Ostern nach dem 22. März eintrifft. So oft also diese Summe nicht
über 9 geht, fällt Ostern noch in den März, und das Datum ergibt sich, wenn
man zu der gedachten Summe 22 addirt. Ist aber jene Summe grösser als 9,
so fällt Ostern in den April, und das Datum erhält man, indem man von der
erwähnten Summe 9 abzieht.

Ich will diese Regeln mit ein Paar Beispielen erläutern. Es sei zuerst
Ostern für das Jahr 1801 zu berechnen. Hier gibt

\textit{nach der ersten Regel} die Division mit 19 den Quotienten 94, um welchen
wir uns nicht bekümmern, und den Rest 15, welches also der erste Rest ist.

\textit{Nach der zweiten Regel} wird der neunzehnfache erste Rest 285; dazu 23
addirt, gibt 308; dies mit 30 dividirt, bleibt der zweite Rest 8.

\textit{Nach der dritten Regel} wird der dritte Rest 1.
\SourceLeaf{89}%

\begin{center}\small
\begin{tabular}{llr}
\textit{Nach der vierten Regel} & addiren wir die vierfache Jahrzahl &7204\\
&den sechsfachen zweiten Rest &48\\
&den doppelten dritten Rest\quad . &2\\
&und die Zahl\quad .\ .\ .\ .\ .\ . &4\\\cline{3-3}
&\hfill gibt Summe &7258
\end{tabular}
\end{center}
Diese mit 7 dividirt, bleibt der vierte Rest 6.

\textit{Nach der fünften Regel} endlich haben wir die Summe des zweiten und vierten
Restes 14; da diese um 5 grösser ist als 9, so fällt Ostern den 5. April. —
Der grüne Donnerstag dieses Jahrs, dessen Andenken jetzt durch die Geschichte
des Tages wieder erneuert wird, fiel also auf den 2. April.

Für das Jahr 1808, welches dem Hrn. Superint. \textsc{Helmuth} zu seinem Aufsatze
Veranlassung gegeben hat, findet sich auf gleiche Weise

nach der ersten Regel der Rest 3.

nach der zweiten Regel das neunzehnfache des ersten Restes 57, dazu
23 gibt 80; hieraus der zweite Rest 20.

nach der dritten Regel der dritte Rest 0.

Nach der vierten Regel haben wir zu addiren 7232, 120, 0 und 4, gibt
die Summe 7356, also der vierte Rest 6.

Nach der fünften Regel erhalten wir die Summe 26; davon also 9 abgezogen,
ergibt sich der 17. April als das Datum des Ostertags.

Nach diesen Erläuterungen wird jeder, dem nur die vier Species bekannt
sind, die Rechnung für andere Jahre ohne Weiteres machen können.

So wie diese Regeln oben gegeben sind, gelten sie für die hundert Jahre
von 1800 bis 1899, für andere Jahrhunderte erleiden sie aber einige Abänderungen.
In dem vergangenen Jahrhundert von 1700 bis 1799 findet weiter kein
Unterschied Statt, als dass in derjenigen Summe, die nach der vierten Regel zu
machen ist, anstatt der Zahl 4 die Zahl 3 genommen werden muss. In dem
siebenzehnten Jahrhundert hingegen, oder vielmehr von 1583 (wo Ostern zum
erstenmal nach der neuen Einrichtung bestimmt wurde) bis 1699 muss nicht nur
in der ebengedachten vierten Regel anstatt der Zahl 4 die Zahl 2 gebraucht,
sondern noch ausserdem in der zweiten Regel zu dem 19fachen ersten Reste nicht
23 sondern 22 hinzugefügt werden. Alles übrige bleibt ungeändert. Blos eben
diese beiden Punkte betreffen nun auch die Abänderungen, die in den nächstkünftigen
Jahrhunderten zu machen sind. Nemlich statt der Zahl 23, die in der
\SourceLeaf{90}%
ersten Regel für das achtzehnte und neunzehnte Jahrhundert gilt, hat man im
zwanzigsten, ein und zwanzigsten und zwei und zwanzigsten die Zahl 24, im
drei und zwanzigsten 25, im vier und zwanzigsten 26, im fünf und zwanzigsten
wieder 25 zu setzen; und statt der Zahl 4, die bei der vierten Regel im neunzehnten,
oder 3, welche im achtzehnten Jahrhundert gebraucht wird, muss
man nehmen
\begin{center}\small
\begin{tabular}{lr}
im zwanzigsten und ein und zwanzigsten\quad .\ .\ .\ . &5\\
im zwei und zwanzigsten\quad .\ .\ .\ .\ .\ .\ .\ .\ . &6\\
im drei und zwanzigsten\quad .\ .\ .\ .\ .\ .\ .\ .\ .\ . &0\\
im vier und zwanzigsten und fünf und zwanzigsten &1
\end{tabular}
\end{center}
Die allgemeinen Regeln für diese Abänderungen findet man am angef. Orte
der Mon. Corresp. [S. 78. d. B.]

Ich habe nur noch zu bemerken, dass die Einrichtung der Epactentafel im
Gregorianischen Kalender so gemacht ist, dass in einzelnen Jahren von den obigen
Regeln zweierlei Ausnahmen Statt finden können. Die erste besteht darin,
dass allemal, wo der zweite Rest 29 und der vierte Rest 6 wird, und wo also
nach der fünften Regel Ostern den 26. April. sein sollte, dieses Fest eine Woche
früher gefeiert wird, also den 19. April. Dieser Fall ist aber bisher nur einmal
eingetreten, nemlich im Jahre 1609, und wird sich nicht eher wieder begeben
als im Jahr 1981. Die zweite Ausnahme kann Statt finden, wenn der zweite
Rest 28 und der vierte 6 wird, wo also nach der fünften Regel Ostern am $25^{\mathrm{sten}}$
April sein sollte, aber wirklich auf den 18. April versetzt wird, doch nur dann
wenn der erste Rest nicht unter 11 war, in welchem Fall die Ausnahme nicht
Statt findet, sondern Ostern auf dem 25. April bleibt. Diese zweite Ausnahme
ist aber bisher noch nie eingetreten, und wird zum erstenmale erst im Jahr 1954
vorkommen.

Alles Vorhergehende betrifft die Berechnung des Osterfestes im Gregorianischen
Kalender. Im Julianischen Kalender bedürfen die obigen fünf Regeln
weiter keiner Abänderung, als dass in der zweiten Regel statt der Zahl 23 die
Zahl 15, und in der vierten Regel statt 4 die Zahl 6 gebraucht werden muss.
Mit diesen Änderungen gelten die Regeln denn allgemein für alle Jahrhunderte.
Auch diese Rechnung will ich mit einem Beispiel erläutern. Für das Jahr 1808
hatten wir nach der ersten Regel, welche ungeändert bleibt, den ersten Rest 3.
Dieser mit 19 multiplicirt und 15 addirt gibt 72, also mit 30 dividirt, wird
\SourceLeaf{91}%
der zweite Rest 12. Der dritte Rest ist wie oben 0. Nach der vierten Regel
haben wir nun zu addiren 7232, 72, 0 und 6, welches gibt 7310, also der vierte
Rest 2. Dessen Summe mit dem zweiten Rest wird 14, wovon nach der fünften
Regel 9 abgezogen der Rest 5 bleibt. Ostern ist also den 5. April nach dem
Julianischen Styl, welches der 17. April nach dem Gregorianischen ist. In diesem
Jahr fällt also Ostern zufälligerweise in beiden Kalendern auf einerlei Tag,
obwohl der Julianische ein 12 Tage kleineres Datum zählt.

Es lässt sich übrigens zeigen, dass im 143. Jahrhundert, d. i. von 14200
bis 14299 Ostern in beiden Kalendern immer auf einerlei Datum fallen würde,
wenn dieselben so lange ungeändert im Gebrauch bleiben sollten, was freilich
nicht zu erwarten ist. Einerlei Datum würde aber alsdann in den beiden Kalendern
immer um 15 volle Wochen auseinander sein, und also die Julianischen
Ostern um 15 Wochen später eintreten als die Gregorianischen. Mit dem Jahre
14300 hört diese Übereinstimmung wieder auf, um 43200 von neuem anzufangen,
wo aber der Unterschied 46 Wochen betragen würde.

Schliesslich bemerke ich, dass im Maiheft der Mon. Corresp. 1802 [S. 80.
d. B.] auch eine rein arithmetische Regel für die Berechnung des jüdischen Osterfestes
gegeben ist, deren Anwendung von aller weitern Bekanntschaft mit der
Einrichtung des jüdischen Kalenders unabhängig ist.

\begin{center}\rule{.25\linewidth}{.3pt}\end{center}
